\documentclass[11pt]{article}

\usepackage[margin=1in]{geometry}
\usepackage{amsmath,amssymb,amsthm,mathtools}
\usepackage{enumitem}
\usepackage{booktabs}
\usepackage{tabularx}
\usepackage{microtype}
\usepackage[hidelinks]{hyperref}
\usepackage{xurl}

\newtheorem{definition}{Definition}
\newtheorem{lemma}{Lemma}
\newtheorem{remark}{Remark}

\title{Sentence Locks and Inline Provenance for Successor Maintenance:\\[0.5ex]
\large Pair-Anchored Reuse, Recheck Guards, and Visible Capsules}
\author{}
\date{Unified archive note v0.09 (2026-03-28; archive v733)}

\begin{document}
\maketitle

\begin{abstract}
A downstream clause pack answers which short checked sentence is available.
That is not yet enough to say whether the sentence may still be reused verbatim for a later cut.
This note isolates the missing policy layer.
It defines a \emph{sentence lock} that anchors one emitted sentence to a concrete base$\to$successor pair, the compare-report cut, the exact shipped object digests that make the sentence true, the drifts that force recheck versus regeneration, the tiny witness fields that may be refreshed in place, the explicit maintained reuse posture that tells later notes whether the wording is still carryforward-clean or only refreshable after recheck, and the small visible-provenance capsule that may travel beside unchanged wording.
The result is a reusable-sentence discipline that is terse, auditable, and separate from the wording note itself.
\end{abstract}

\paragraph{Non-redundancy note.}
Synthesis~31 owns the downstream wording suffix---summary, support, citation, quote, and clause.\cite{series:downstreamnormalforms}
Synthesis~32 owns the stable release spine.\cite{series:releasespine}
Synthesis~34 owns the status-envelope / lineage-notice split at the summary stage.\cite{series:statusenvelopes}
Synthesis~46 owns whole-question answer sheets, Synthesis~47 owns emitted answer cards built from those sheets plus one sentence lock, Synthesis~51 owns publication profiles that pair one emitted answer with one inherited publication budget and one escalation path, Synthesis~54 owns the answer-side terminal citation normal form, Synthesis~74 owns the paired terminal citation discipline, and Synthesis~75 owns the paired cutover rule.\cite{series:challengeanswersheets,series:challengeanswercards,series:challengeanswerpublicationprofiles,series:challengeanswerreviewmenus,series:pairedterminalcitationdiscipline,series:pairedterminalcutoverrules}
This note owns only the policy for reusing an already-checked emitted sentence without silently letting it drift across release pairs or object cuts. Later notes should reopen it only when the exact lock basis, recheck boundary, refresh-in-place boundary, or visible inline-provenance posture is itself live.

\paragraph{Scope.}
This is not a new mechanism note and not a new release stage.
It is a publication-interface note about what must remain fixed before a later paper may print the same checked sentence again.

\paragraph{Imports.}
Contract-object vocabulary is imported from Synthesis~0.\cite{series:contractobjects}
Drift language is imported from Synthesis~18.\cite{series:planstateequiv}
Replay/closure owners are imported from Synthesis~20 and Synthesis~33.\cite{series:auditorreplay,series:releaseclosure}
The downstream wording interface is imported from Synthesis~31.\cite{series:downstreamnormalforms}
Whole-question answer sheets, emitted answer cards, and answer publication profiles are imported from Synthesis~46, Synthesis~47, and Synthesis~51 whenever later notes want the audience/question handoff, the final released answer object, or a budgeted publication package rather than the reuse policy itself.\cite{series:challengeanswersheets,series:challengeanswercards,series:challengeanswerpublicationprofiles}
The answer-side terminal citation normal form is imported from Synthesis~54, one concrete checked-answer/request-tail seam is imported from Synthesis~17, and the abstract paired reuse/cutover discipline is imported from Synthesis~74--75 as the later-paper stop rule around this sentence-lock owner.\cite{series:challengeanswerreviewmenus,series:workedexample,series:pairedterminalcitationdiscipline,series:pairedterminalcutoverrules}
The concrete worked instantiation is imported from Synthesis~17.\cite{series:workedexample}

\paragraph{Exports.}
This note exports (i) sentence locks, (ii) a recheck-vs-regenerate split, (iii) an in-place refresh boundary, (iv) a visible inline-provenance profile discipline, and (v) a local-reopen rule saying that later terse papers should reopen this note only when the exact lock basis, recheck boundary, refresh-in-place boundary, or visible inline-provenance posture is itself live. Otherwise later terse papers should default to the answer-side terminal citation normal form in Synthesis~54, drop to Synthesis~17 when one concrete checked-answer/request-tail seam is wanted, and widen to Synthesis~74--75 only for the abstract paired terminal discipline or abstract stop/widen/reopen judgment. Whole-question handoff, final emitted audience/question answer objects, and budgeted publication packaging remain outside scope here and belong to Synthesis~46--47 and Synthesis~51.\cite{series:challengeanswersheets,series:challengeanswercards,series:challengeanswerpublicationprofiles,series:challengeanswerreviewmenus,series:workedexample,series:pairedterminalcitationdiscipline,series:pairedterminalcutoverrules}

\paragraph{Later-terse-paper citation posture.}
When a later short paper only needs to answer \emph{whether one already-checked successor sentence may still travel unchanged, travel only after recheck, or must be retired}, it should stop at the minimal public sentence-lock / inline-provenance card below plus the tiny lock-forcing drill. It should widen only when the live issue becomes downstream wording ownership (Synthesis~31), summary-stage public-status compression (Synthesis~34), whole-question answer-sheet or answer-card ownership (Synthesis~46--47), publication-budget packaging (Synthesis~51), the maintained checked-answer handoff itself (Synthesis~54), one concrete checked-answer/request-tail seam (Synthesis~17), or the abstract paired-tail discipline / abstract stop-widen-reopen rule (Synthesis~74--75).\cite{series:downstreamnormalforms,series:statusenvelopes,series:challengeanswersheets,series:challengeanswercards,series:challengeanswerpublicationprofiles,series:challengeanswerreviewmenus,series:workedexample,series:pairedterminalcitationdiscipline,series:pairedterminalcutoverrules}

\paragraph{One branched pre-answer route.}
When a later terse paper challenges one already reused checked sentence, the honest owner map is now: Synthesis~35 for whether that visible sentence was allowed to travel unchanged at all and what inline provenance floor traveled with it; Synthesis~36 for which audience/question pair should start where; Synthesis~37 for which semantic piece and approved terminal field family are even in scope; and Synthesis~38 for the exact contested-piece walk. This note is therefore an optional sidecar, not the challenge-route root. Stop here only while the live issue is carryforward / recheck / retire posture or inline provenance mode for the challenged sentence itself. Hand off to Synthesis~36--38 as soon as the live issue becomes audience entry, semantic stop semantics, or exact contested-piece escalation.\cite{series:challengeroutes,series:challengestops,series:challengebranches}

\paragraph{One tiny sidecar drill.}
If a later terse paper asks whether the already-shipped public-status sentence may still be reused unchanged, Synthesis~35 is the honest first stop. Once the next live question becomes where a referee should begin, which continuity piece may stop at which verdict, or what exact path the continuity challenge takes, the honest owners become Synthesis~36, Synthesis~37, and Synthesis~38 respectively; quoting the sentence-lock note for those later questions would already be one layer off.\cite{series:challengeroutes,series:challengestops,series:challengebranches}

\section{One sentence, one lock}

\begin{definition}[Sentence lock]
For one emitted downstream sentence $\sigma$, a \emph{sentence lock}
\[
L(\sigma)=(r_{\mathrm{base}},r_{\mathrm{succ}},d,H,T)
\]
consists of one base release id, one successor release id, one compare-report id, one finite digest witness $H$ for the concrete shipped objects that make $\sigma$ true, and one finite terminal-target set $T$ naming the narrow fields whose values are compressed by $\sigma$.
A sentence is reusable only while this lock remains valid for the later reuse site.
\end{definition}

\begin{definition}[Recheck vs regeneration]
For one sentence lock $L(\sigma)$, a \emph{recheck set}
\[
R_{\mathrm{chk}}(\sigma)
\]
contains drifts that require the sentence to be revalidated against a new shipped object cut before verbatim reuse.
A \emph{regeneration set}
\[
R_{\mathrm{regen}}(\sigma)
\]
contains drifts that retire the wording and force a newly emitted sentence.
Typical recheck drifts are object-digest changes under the same release pair and same terminal targets; typical regeneration drifts are release-pair changes or terminal-target changes.
\end{definition}

\begin{definition}[Maintained reuse posture]
For one sentence lock $L(\sigma)$, a \emph{maintained reuse posture}
\[
S(\sigma)\in\{\texttt{carryforward},\texttt{refresh\_only\_after\_recheck},\texttt{retire\_and\_reissue}\}
\]
is the current maintained status exported beside the lock. The first value means the already-checked sentence may presently travel unchanged under its default lock conditions; the second means the wording is still the governing sentence family but may only be carried forward after the declared recheck guards are discharged for the current cut; the third means the wording is no longer reusable and a new emitted sentence is required.
\end{definition}

\begin{definition}[Refresh-in-place boundary]
For one sentence lock $L(\sigma)$, an \emph{in-place refresh set}
\[
F_{\mathrm{ref}}(\sigma)
\]
is the finite list of witness fields that may be updated after every triggered recheck guard has been discharged while every regeneration guard remains false.
It excludes the emitted text, the release-pair anchor, and the terminal targets that make the sentence true.
\end{definition}

\begin{remark}[Why this is separate from the normal form]
The downstream normal form answers what a sentence resolves to.
A sentence lock answers whether the already-checked sentence may still be reused for the present cut.
Those are different questions and they should not be merged.
\end{remark}

\begin{remark}[Why posture helps]
Later terse notes often need one bit of maintenance policy more than they need the full lock internals: may this exact checked sentence still travel, may it travel only after a fresh recheck, or must it be retired. Exporting that posture directly avoids forcing each neighboring paper to reconstruct the answer from lock guards and digest clauses.
\end{remark}

\section{Visible provenance is small on purpose}

\begin{definition}[Witness capsule and witness floor]
For one sentence lock $L(\sigma)$, a \emph{witness capsule} $W(\sigma)$ is the finite set of already-declared anchor fields that may travel beside unchanged wording as visible provenance.
A \emph{witness floor} $F_{\mathrm{vis}}(\sigma)\subseteq W(\sigma)$ is the smallest declared subset that should appear whenever any inline provenance is shown at all.
\end{definition}

\begin{definition}[Inline provenance profile]
For one sentence lock with witness capsule $W(\sigma)$ and witness floor $F_{\mathrm{vis}}(\sigma)$, an \emph{inline provenance profile} is one of three declared modes:
\[
P_{\mathrm{silent}}(\sigma)=\varnothing,
\qquad P_{\mathrm{floor}}(\sigma)=F_{\mathrm{vis}}(\sigma),
\qquad P_{\mathrm{capsule}}(\sigma)=W(\sigma).
\]
A \emph{default inline provenance profile} is the single declared mode later notes should inherit unless they explicitly choose another.
\end{definition}

\begin{remark}[Why this stays narrow]
The provenance rule does not enlarge what may travel beside a sentence.
It only states whether visible provenance is absent, reduced to the minimum declared bundle, or expanded to the full declared capsule.
That keeps terse reuse sites from improvising footnote policy sentence by sentence.
\end{remark}

\section{Reuse rule}

\begin{lemma}[Locked reuse discipline]
Let $\sigma$ be an emitted sentence with a valid sentence lock $L(\sigma)$, a valid recheck set, a valid regeneration set, and a valid in-place refresh boundary.
Then verbatim reuse of $\sigma$ is auditable: if a regeneration guard fires, new wording is required; if only recheck guards fire, the sentence may be reused only after revalidation and only the declared refresh-in-place witness fields may be updated; if no guard fires, the sentence may be reused under its declared inline-provenance policy.
\end{lemma}

\begin{proof}
The sentence lock fixes the release pair, compare-report cut, supporting object digests, and terminal targets that make the sentence true.
The regeneration set marks drifts that change the comparison itself or the narrow targets the sentence compresses, so verbatim reuse would no longer be honest.
The recheck set marks smaller drifts where the sentence may remain true but only after validation against the new shipped object cut.
The refresh boundary then prevents maintainers from mutating the sentence or its truth-bearing targets while pretending that only provenance was updated.
Thus locked reuse preserves the distinction between checked wording and later maintenance housekeeping.
\end{proof}

\begin{table}[t]
\centering
\small
\begin{tabularx}{\linewidth}{@{}l l X@{}}
\toprule
Layer & Canonical object & Question answered \\
\midrule
Wording suffix & downstream normal form & What does this public sentence resolve to if challenged? \\
Reuse policy & sentence lock & May this exact checked sentence still be reused for this pair and object cut? \\
Inline display & provenance profile & If the sentence is reused, how much visible provenance should travel beside it? \\
\bottomrule
\end{tabularx}
\caption{The wording note and the reuse-policy note answer different questions and stay cleaner when they are separated.}
\label{tab:locksplit}
\end{table}

\paragraph{A minimal public sentence-lock / inline-provenance card.}
\begin{table}[t]
\centering
\small
\begin{tabularx}{0.97\linewidth}{@{}p{0.31\linewidth}X@{}}
\toprule
\textbf{Public row} & \textbf{Pinned value}\\
\midrule
Release-pair anchor & \texttt{base\_release\_id = worked-example-draft-194}, \texttt{successor\_release\_id = worked-example-draft-194b}, and \texttt{compare\_report\_id = worked-example-compare-report-v1} for the worked public-status sentence \\
Checked sentence family & The maintained sentence family is the worked lineage sentence saying that the successor remains on the same certified interface after replaying changed fallback surfaces and that required successor publication outputs are complete \\
Digest / target floor & The sentence lock binds the exact shipped object digests that make that sentence true together with the terminal-target set covering the status envelope and lineage-notice fields compressed by the sentence \\
Guard split & Release-pair or terminal-target drift belongs to \texttt{regenerate}; same-pair shipped-object digest drift under the same targets belongs to \texttt{recheck} \\
Maintained reuse posture & The worked lock currently exports a declared reuse posture together with a finite in-place refresh set for witness fields only \\
Visible provenance mode & The default inline-provenance profile is inherited from the worked lock and may expose either silence, the witness floor, or the full witness capsule without changing the checked sentence \\
Escalation pointer & Widen to Synthesis~31 for wording ownership, Synthesis~34 for summary-stage wrapper ownership, Synthesis~46--47 for whole-question answer objects, Synthesis~51 for publication packaging, stop at Synthesis~54 for the maintained checked-answer handoff, drop to Synthesis~17 for one concrete checked-answer/request-tail seam, and widen to Synthesis~74--75 only for the abstract paired-tail discipline or abstract stop/widen/reopen judgment once the live issue is no longer sentence reuse itself \\
\bottomrule
\end{tabularx}
\caption{A minimal public sentence-lock / inline-provenance card. This is the smallest public answer later terse papers can quote directly when the live issue is whether one checked sentence still travels unchanged, travels only after recheck, or must be retired.}
\label{tab:minimal-sentence-lock-card}
\end{table}

This card is intentionally non-decorative: it fixes one release pair, one checked sentence family, one digest-and-target floor, one recheck-vs-regenerate split, one maintained reuse posture, and one visible-provenance mode without silently re-owning the wording suffix, the whole-question answer object, or the final publication budget.\cite{series:workedexample}

\paragraph{A forced public sentence-lock / inline-provenance matrix.}
\begin{table}[t]
\centering
\small
\begin{tabularx}{\linewidth}{@{}p{0.22\linewidth}p{0.46\linewidth}X@{}}
\toprule
\textbf{Shortcut} & \textbf{Still the same packet iff} & \textbf{Otherwise honest action}\\
\midrule
Sentence lock & The maintained cut still has the same release-pair anchor, compare-report cut, digest witness family, and terminal-target set for the sentence & Reopen this note or widen to Synthesis~31 / Synthesis~34 instead of citing the old lock \\
Reuse posture & The same sentence lock still holds and the triggered-guard pattern still yields the same declared posture \texttt{carryforward}, \texttt{refresh\_only\_after\_recheck}, or \texttt{retire\_and\_reissue} & Stop at the changed posture and do not reuse the old sentence policy \\
Inline provenance profile & The same sentence lock still holds and the same visible-provenance mode and witness floor / capsule remain declared & Reopen this note rather than improvising a new inline provenance tail \\
\bottomrule
\end{tabularx}
\caption{Forced public sentence-lock / inline-provenance matrix. Sentence reuse, refresh-only rechecks, and visible provenance remain quotable only while the same pair anchor, digest-and-target floor, guard split, and provenance mode still agree.}
\label{tab:forced-sentence-lock-matrix}
\end{table}

\paragraph{A compact sentence-lock first-reopen-owner card.}
\begin{table}[t]
\centering
\small
\begin{tabularx}{\linewidth}{@{}lX@{}}
\toprule
Field & First reopened owner \\
\midrule
Same lock still owns & Reopen Synthesis~35 first only while the same release-pair anchor, compare-report cut, digest-and-target floor, declared recheck-vs-regenerate split, maintained reuse posture, and visible inline-provenance mode still force the same sentence-lock contract, and the live issue is still whether one checked sentence may travel unchanged, only after recheck, or not at all \\
Narrower imported owners & Reopen Synthesis~31 first if the wording suffix or smallest sufficient citation target for the sentence changes, and reopen Synthesis~34 first if the visible status-envelope / lineage-notice wrapper class becomes the honest moving part even though the same checked sentence family would otherwise look stable \\
Immediate right boundary & Reopen Synthesis~36 first if the lock itself still holds but the live issue becomes which audience/question pair should begin where; reopen Synthesis~37 first if that route remains fixed but the semantic stop family or approved terminal field split moves; widen to Synthesis~38 only after those imports are fixed and the exact contested-piece walk becomes the honest moving part rather than sentence reuse itself \\
Later answer-side boundary & Widen to Synthesis~46--47 when the live issue has become the whole-question answer object or the final emitted answer object, widen to Synthesis~51 when the governing issue is publication-budget packaging, stop at Synthesis~54 for the maintained checked-answer handoff, drop to Synthesis~17 for one concrete checked-answer/request-tail seam, and widen to Synthesis~74--75 only for the abstract paired-tail discipline or abstract stop/widen/reopen judgment once sentence reuse is no longer the moving part \\
Selection rule & Later terse papers may quote this card only to recover the first narrower or wider owner after a stale checked-sentence shortcut; they should not restate stable audience-route, stop-profile, exact-branch, answer-card, or paired-terminal owners here once one of those neighbors has become the honest moving part \\
\bottomrule
\end{tabularx}
\caption{A compact sentence-lock first-reopen-owner card. This is the smallest owner-cutover map later terse papers can quote directly when a maintained checked-sentence shortcut has gone stale and the remaining task is to recover which narrower or wider owner regains authority first.}
\label{tab:sentence-lock-first-reopen-owner-card}
\end{table}

\paragraph{A tiny first-reopen-owner drill.}
For the worked public-status lineage sentence, a later terse paper may keep citing the sentence-lock owner only while the same base$\to$successor pair, compare-report cut, digest-and-target floor, guard split, and inline-provenance mode still force the same checked-sentence reuse contract. Reopen Synthesis~34 first if the live issue becomes whether the outward-facing lineage-notice wrapper is still the honest public face, reopen Synthesis~36 first if the lock itself still holds but the referee or release-note entry route changes, reopen Synthesis~37 first if that route remains fixed but the continuity-versus-closure stop semantics move, and widen to Synthesis~38 only after those imports are fixed and the exact continuity-piece walk becomes the honest moving part. Widen to Synthesis~46--47 or Synthesis~51 only after the sentence-reuse contract itself is no longer the honest point of disagreement; then stop at Synthesis~54 for the maintained checked-answer handoff, drop to Synthesis~17 for one concrete checked-answer/request-tail seam, or widen to Synthesis~74--75 only when the remaining issue is the abstract paired-tail discipline or abstract stop/widen/reopen judgment.\cite{series:workedexample,series:downstreamnormalforms,series:statusenvelopes,series:challengeroutes,series:challengestops,series:challengebranches,series:challengeanswersheets,series:challengeanswercards,series:challengeanswerpublicationprofiles,series:challengeanswerreviewmenus,series:pairedterminalcitationdiscipline,series:pairedterminalcutoverrules}

\paragraph{One tiny lock-forcing drill.}
In the worked Case-B successor, later terse papers may keep citing the same checked lineage sentence only while the release-pair anchor stays fixed, the sentence still compresses the same status-envelope / lineage-notice targets, and no regeneration guard has fired. If a later maintained cut changed only shipped-object digests under the same pair and same targets, then the honest posture would move to \texttt{refresh\_only\_after\_recheck} rather than free carryforward. If the release pair or terminal targets changed, then the old lock would already be stale and the sentence would need to be retired and reissued here before any later paper could honestly reuse it.\cite{series:workedexample}

\section{Worked example pointer}
The maintained worked bundle now ships \path{example_successor_sentence_locks.json}.\cite{series:workedexample}
It binds two emitted sentences for the canned Case-B successor: the outward-facing public-status sentence and the compact numeric diff sentence.
For each one, the object records the release-pair anchor, compare-report cut, terminal targets, exact shipped object digests, recheck-vs-regenerate guards, refresh-in-place witness fields, witness capsule, witness floor, optional visible-provenance tail, and default inline provenance profile.
The companion \path{example_successor_normal_form.json} continues to own the suffix itself, while the newer \path{example_successor_challenge_answer_cards.json} imports one sentence lock together with one whole-question answer sheet whenever a later note wants the final emitted audience/question answer object rather than the reuse policy alone. The newer \path{example_successor_challenge_answer_publication_profiles.json} then packages those same sentence-lock modes for brief publication without becoming the owner of sentence retirement. The maintained routing/bridge adjuncts, \path{example_question_routes.json} and \path{example_series_spine.json}, now point directly at that answer-card / sentence-lock pair together with the default inline-provenance mode, so later terse papers can inherit one checked sentence and one declared visible-provenance floor without re-deriving those bindings from the answer stack by hand. The sentence-lock ledger exists so later papers can cite reuse policy without treating the normal-form note as a sentence-retirement manual, the answer-card note as the owner of sentence retirement, or the publication-profile note as the owner of visible-provenance semantics.\cite{series:downstreamnormalforms,series:challengeanswercards,series:challengeanswerpublicationprofiles,series:seriesspines}

\paragraph{Why this helps the series.}
This split lets Synthesis~31 stay about wording and lets the release-spine notes stay about stage ownership.
The worked example remains auditable, but later papers no longer need to pull sentence-retirement policy out of the middle of the wording note.

\begin{thebibliography}{99}\small

\bibitem{series:contractobjects}
Anonymous.
\newblock Contract Objects for Anonymous-DHT Claims: Commitments, Menus, Receipts, and Release Bindings.
\newblock Series synthesis note (Synthesis~0), 2026. (This archive.)

\bibitem{series:planstateequiv}
Anonymous.
\newblock Plan and State Equivalence: Digest-Bound Replay Hooks and Drift Taxonomy.
\newblock Series synthesis note (Synthesis~18), 2026. (This archive.)

\bibitem{series:auditorreplay}
Anonymous.
\newblock Auditor Replay Recipe for Anonymous-DHT Receipts: A Minimal Checklist and Verdict Tuple.
\newblock Series synthesis note (Synthesis~20), 2026. (This archive.)

\bibitem{series:releaseclosure}
Anonymous.
\newblock Release Obligations and Closure Ledgers for Successor Maintenance: Required Roles, Row Satisfaction, and Completion Tokens.
\newblock Series synthesis note (Synthesis~33), 2026. (This archive.)

\bibitem{series:downstreamnormalforms}
Anonymous.
\newblock Downstream Normal Forms for Successor Maintenance: Summary, Support, Citation, Quote, and Clause Discipline.
\newblock Series synthesis note (Synthesis~31), 2026. (This archive.)

\bibitem{series:releasespine}
Anonymous.
\newblock Stable Release Spine Contracts for Successor Maintenance: Stage Ownership, Handoff Rules, and Auditable Walk-Throughs.
\newblock Series synthesis note (Synthesis~32), 2026. (This archive.)

\bibitem{series:statusenvelopes}
Anonymous.
\newblock Successor Status Envelopes and Lineage Notices for Release Maintenance: Narrow Status Tokens, Reader Pointer Order, and Public Answers.
\newblock Series synthesis note (Synthesis~34), 2026. (This archive.)

\bibitem{series:challengeroutes}
Anonymous.
\newblock Challenge Routes and Audience Profiles for Successor Maintenance: Audience-Keyed Start Objects and Escalation Order.
\newblock Series synthesis note (Synthesis~36), 2026. (This archive.)

\bibitem{series:challengestops}
Anonymous.
\newblock Challenge Stop Profiles for Successor Maintenance: Question-Keyed Terminal Fields and Semantic Completion Rules.
\newblock Series synthesis note (Synthesis~37), 2026. (This archive.)

\bibitem{series:challengebranches}
Anonymous.
\newblock Challenge Branch Maps for Successor Maintenance: Piece-Keyed Escalation Walks and Exact Terminal Owners.
\newblock Series synthesis note (Synthesis~38), 2026. (This archive.)

\bibitem{series:challengeanswersheets}
Anonymous.
\newblock Challenge Answer Sheets for Successor Maintenance: Audience-Question Piece-Complete Handoffs from Answer Capsules.
\newblock Series synthesis note (Synthesis~46), 2026. (This archive.)

\bibitem{series:challengeanswercards}
Anonymous.
\newblock Challenge Answer Cards for Successor Maintenance: Audience-Question Released Answers from Answer Sheets and Sentence Locks.
\newblock Series synthesis note (Synthesis~47), 2026. (This archive.)

\bibitem{series:challengeanswerpublicationprofiles}
Anonymous.
\newblock Challenge Answer Publication Profiles for Successor Maintenance: Budgeted Reuse Modes from Answer Cards, Sentence Locks, and Citation Dockets.
\newblock Series synthesis note (Synthesis~51), 2026. (This archive.)


\bibitem{series:seriesspines}
Anonymous.
\newblock Series Spines from Mathematical Roots to Referee Handoffs: A Non-Redundant Bridge from Backbones, Receipts, Replay, Release Evolution, and Review Menus.
\newblock Series synthesis note (Synthesis~55), 2026. (This archive.)

\bibitem{series:challengeanswerreviewmenus}
Anonymous.
\newblock Challenge Answer Review Menus for Successor Maintenance: Answer-Side Terminal Citation Normal Form.
\newblock Series synthesis note (Synthesis~54), 2026. (This archive.)

\bibitem{series:pairedterminalcitationdiscipline}
Anonymous.
\newblock Paired Terminal Citation Discipline for Review Spines: Stable Answer Stops and Adjacent Request Widening.
\newblock Series synthesis note (Synthesis~74), 2026. (This archive.)

\bibitem{series:pairedterminalcutoverrules}
Anonymous.
\newblock Paired Terminal Cutover Rules for Review Spines: When Later Terse Papers Stop, Widen, or Reopen.
\newblock Series synthesis note (Synthesis~75), 2026. (This archive.)

\bibitem{series:workedexample}
Anonymous.
\newblock Worked Example for Anonymous-DHT Receipts: Minimal Public Surface, Cross-Release Diff, and Successor Interlock.
\newblock Series synthesis note (Synthesis~17), 2026. (This archive.)

\end{thebibliography}
\end{document}
